\documentclass[11pt]{article}
\usepackage[utf8]{inputenc}
\usepackage{amsmath, amssymb, amsthm}
\usepackage{algorithm}
\usepackage{algpseudocode}
\usepackage{geometry}
\usepackage{hyperref}
\usepackage{graphicx}
\usepackage{booktabs}
\usepackage{listings}
\usepackage{xcolor}

\geometry{a4paper, margin=1in}

\lstset{
    basicstyle=\ttfamily\small,
    breaklines=true,
    frame=single,
    language=Python,
    keywordstyle=\color{blue},
    commentstyle=\color{gray},
    stringstyle=\color{red}
}

\title{CMPE 300: Analysis of Algorithms\\Project 3 - Interactive Proof Systems\\Final Report}
\author{Ali Ayhan Gunder 2021400219 \and Yigit Memceoktay 2020403076}
\date{\today}

\begin{document}

\maketitle

\tableofcontents
\newpage

%==============================================================================
\section{Introduction}
%==============================================================================

This report presents our solutions and analysis for Project 3, which explores Interactive Proof Systems in the context of computational complexity theory. The project covers zero-knowledge proofs for graph 3-coloring and graph isomorphism problems, along with a practical demonstration of exploiting a biased verifier in the 3-coloring protocol.

Interactive proof systems are fundamental to modern cryptography and complexity theory, allowing a \textit{prover} to convince a \textit{verifier} of a statement's truth without revealing unnecessary information. The key properties we analyze are:
\begin{itemize}
    \item \textbf{Completeness}: An honest prover can always convince the verifier.
    \item \textbf{Soundness}: A dishonest prover cannot fool the verifier (except with negligible probability).
    \item \textbf{Zero-Knowledge}: The verifier learns nothing beyond the truth of the statement.
\end{itemize}

%==============================================================================
\section{Q1: Interactive Algorithm for 3-Coloring}
%==============================================================================

\subsection{Part (a): Problem Definition}

A graph $G = (V, E)$ is \textbf{3-colorable} if there exists a function $c: V \to \{0, 1, 2\}$ such that for every edge $(u, v) \in E$, we have $c(u) \neq c(v)$. In other words, no two adjacent vertices share the same color.

The corresponding decision problem is defined as:
\[
\textsc{3Color} = \{ G \mid G \text{ is a graph that admits a valid 3-coloring} \}
\]

This problem is known to be NP-complete, meaning:
\begin{enumerate}
    \item Given a proposed coloring, we can verify its validity in polynomial time (NP).
    \item Any problem in NP can be reduced to 3-Coloring in polynomial time (NP-hard).
\end{enumerate}

\subsection{Part (b): Zero-Knowledge Interactive Proof Protocol}

\textbf{Setup:}
\begin{itemize}
    \item \textbf{Common Input:} A graph $G = (V, E)$ with $n$ vertices and $m$ edges.
    \item \textbf{Prover's Secret:} A valid 3-coloring $c: V \to \{0, 1, 2\}$.
\end{itemize}

\textbf{Protocol (Single Round):}
\begin{enumerate}
    \item \textbf{Commitment Phase:}
    \begin{itemize}
        \item The Prover selects a random permutation $\pi$ of the colors $\{0, 1, 2\}$.
        \item The Prover computes a permuted coloring $c'(v) = \pi(c(v))$ for all $v \in V$.
        \item For each vertex $v$, the Prover creates a cryptographic commitment $\text{Commit}(c'(v), r_v)$ where $r_v$ is a random string.
        \item The Prover sends all commitments to the Verifier.
    \end{itemize}
    
    \item \textbf{Challenge Phase:}
    \begin{itemize}
        \item The Verifier selects an edge $(u, v) \in E$ uniformly at random.
        \item The Verifier sends $(u, v)$ to the Prover.
    \end{itemize}
    
    \item \textbf{Response Phase:}
    \begin{itemize}
        \item The Prover opens the commitments for vertices $u$ and $v$ by revealing $(c'(u), r_u)$ and $(c'(v), r_v)$.
    \end{itemize}
    
    \item \textbf{Verification Phase:}
    \begin{itemize}
        \item The Verifier checks that the opened values are consistent with the commitments.
        \item The Verifier checks that $c'(u) \neq c'(v)$.
        \item If both checks pass, accept this round; otherwise, reject.
    \end{itemize}
\end{enumerate}

The protocol is repeated $k$ times (e.g., $k = m \cdot \log(1/\epsilon)$ for error bound $\epsilon$) with fresh randomness each round.

\subsection{Part (c): Probabilistic Analysis}

\subsubsection{Completeness}
\textbf{Claim:} If $G$ is 3-colorable and the Prover is honest, the Verifier accepts with probability 1.

\textbf{Proof:} Let $c$ be the valid 3-coloring known to the Prover. For any permutation $\pi$, the permuted coloring $c' = \pi \circ c$ is also valid because permutations are bijections. Specifically:
\begin{itemize}
    \item For any edge $(u, v) \in E$: $c(u) \neq c(v)$ (by validity of $c$)
    \item Since $\pi$ is a bijection: $\pi(c(u)) \neq \pi(c(v))$
    \item Therefore: $c'(u) \neq c'(v)$
\end{itemize}

Thus, for any edge the Verifier challenges, the revealed colors will be different, and the Verifier accepts with probability 1 in each round, and hence probability 1 overall.

\subsubsection{Soundness}
\textbf{Claim:} If $G$ is not 3-colorable, any cheating Prover convinces the Verifier with probability at most $(1 - 1/m)^k$ after $k$ rounds.

\textbf{Proof:} If $G$ is not 3-colorable, then for any coloring $c': V \to \{0, 1, 2\}$ that the Prover commits to, there exists at least one ``bad'' edge $(u, v) \in E$ such that $c'(u) = c'(v)$.

\begin{itemize}
    \item The commitment scheme is binding, so the Prover cannot change colors after committing.
    \item The Verifier selects edges uniformly at random, so $\Pr[\text{select bad edge}] \geq \frac{1}{m}$.
    \item The probability of not catching the Prover in one round is at most $1 - \frac{1}{m}$.
    \item After $k$ independent rounds, the probability of never catching the Prover is at most:
    \[
    \left(1 - \frac{1}{m}\right)^k
    \]
\end{itemize}

For $k = m \ln(1/\epsilon)$, we have:
\[
\left(1 - \frac{1}{m}\right)^{m \ln(1/\epsilon)} \approx e^{-\ln(1/\epsilon)} = \epsilon
\]

Thus, soundness error decreases exponentially with the number of rounds.

\subsubsection{Zero-Knowledge (Intuition)}

The protocol is zero-knowledge because the Verifier learns only the colors of two adjacent vertices per round, and these colors are uniformly distributed due to the random permutation $\pi$.

\textbf{Simulator Construction:}
\begin{enumerate}
    \item The simulator picks two random distinct colors $(a, b)$ from $\{0, 1, 2\}$.
    \item For the challenged edge $(u, v)$, it outputs $(a, b)$.
    \item This distribution is identical to the real protocol since $\pi$ is uniform.
\end{enumerate}

The key insight is that the random permutation $\pi$ ``masks'' the original coloring structure. Each round reveals only a random pair of distinct colors, which could correspond to any valid coloring. The Verifier cannot correlate information across rounds because $\pi$ is chosen fresh each time.

%==============================================================================
\section{Q2: Interactive Algorithm for Graph Isomorphism}
%==============================================================================

\subsection{Part (a): Problem Definition}

Two graphs $G_1 = (V_1, E_1)$ and $G_2 = (V_2, E_2)$ are \textbf{isomorphic} (denoted $G_1 \cong G_2$) if there exists a bijection $\phi: V_1 \to V_2$ such that:
\[
(u, v) \in E_1 \iff (\phi(u), \phi(v)) \in E_2
\]

The decision problem is:
\[
\textsc{GI} = \{ (G_1, G_2) \mid G_1 \cong G_2 \}
\]

Graph Isomorphism is in NP (the isomorphism serves as a polynomial-time verifiable certificate) but is not known to be NP-complete or in P.

\subsection{Part (b): Zero-Knowledge Interactive Protocol}

\textbf{Setup:}
\begin{itemize}
    \item \textbf{Common Input:} Two graphs $G_1, G_2$ on $n$ vertices.
    \item \textbf{Prover's Secret:} An isomorphism $\pi: V_1 \to V_2$ such that $\pi(G_1) = G_2$.
\end{itemize}

\textbf{Protocol (Single Round):}
\begin{enumerate}
    \item \textbf{Commitment:}
    \begin{itemize}
        \item The Prover selects a random permutation $\sigma$ of the $n$ vertices.
        \item The Prover computes $H = \sigma(G_1)$ (equivalently, $H = \sigma \circ \pi^{-1}(G_2)$).
        \item The Prover sends $H$ to the Verifier (or commits to its adjacency matrix).
    \end{itemize}
    
    \item \textbf{Challenge:}
    \begin{itemize}
        \item The Verifier selects a random bit $b \in \{1, 2\}$ uniformly at random.
        \item The Verifier sends $b$ to the Prover.
    \end{itemize}
    
    \item \textbf{Response:}
    \begin{itemize}
        \item If $b = 1$: The Prover sends $\rho = \sigma$ (showing $H \cong G_1$).
        \item If $b = 2$: The Prover sends $\rho = \sigma \circ \pi^{-1}$ (showing $H \cong G_2$).
    \end{itemize}
    
    \item \textbf{Verification:}
    \begin{itemize}
        \item The Verifier checks that $\rho(G_b) = H$.
        \item If the check passes, accept; otherwise, reject.
    \end{itemize}
\end{enumerate}

\subsection{Part (c): Probabilistic Analysis}

\subsubsection{Completeness}
\textbf{Claim:} If $G_1 \cong G_2$ and the Prover knows the isomorphism $\pi$, the Verifier accepts with probability 1.

\textbf{Proof:}
\begin{itemize}
    \item If $b = 1$: Prover sends $\sigma$. Since $H = \sigma(G_1)$, the Verifier confirms $\sigma(G_1) = H$. $\checkmark$
    \item If $b = 2$: Prover sends $\sigma \circ \pi^{-1}$. Since $\pi(G_1) = G_2$ and $H = \sigma(G_1)$:
    \[
    (\sigma \circ \pi^{-1})(G_2) = \sigma(\pi^{-1}(G_2)) = \sigma(G_1) = H \quad \checkmark
    \]
\end{itemize}

\subsubsection{Soundness}
\textbf{Claim:} If $G_1 \not\cong G_2$, any cheating Prover succeeds with probability at most $(1/2)^k$ after $k$ rounds.

\textbf{Proof:} If $G_1 \not\cong G_2$, then any graph $H$ can be isomorphic to at most one of $G_1$ or $G_2$ (but not both).

\begin{itemize}
    \item The Prover must commit to $H$ before knowing $b$.
    \item If $H \cong G_1$ but $H \not\cong G_2$: Prover can only answer $b = 1$ correctly.
    \item If $H \cong G_2$ but $H \not\cong G_1$: Prover can only answer $b = 2$ correctly.
    \item If $H$ is isomorphic to neither: Prover cannot answer either challenge.
\end{itemize}

Thus, the Prover succeeds with probability at most $1/2$ per round. After $k$ rounds:
\[
\Pr[\text{Prover succeeds in all rounds}] \leq \left(\frac{1}{2}\right)^k
\]

\subsubsection{Zero-Knowledge (Intuition)}

The Verifier sees:
\begin{enumerate}
    \item A random graph $H$ isomorphic to both $G_1$ and $G_2$ (since they are isomorphic to each other).
    \item An isomorphism from $G_b$ to $H$.
\end{enumerate}

Since $\sigma$ is chosen uniformly at random:
\begin{itemize}
    \item $H$ is a uniformly random graph from the isomorphism class of $G_1$ (and $G_2$).
    \item The revealed isomorphism is uniformly random among all isomorphisms from $G_b$ to $H$.
\end{itemize}

A simulator can perfectly simulate the transcript by:
\begin{enumerate}
    \item Picking a random $b' \in \{1, 2\}$.
    \item Generating $H = \rho(G_{b'})$ for a random permutation $\rho$.
    \item If $b = b'$, output $\rho$; otherwise, rewind and try again.
\end{enumerate}

The specific isomorphism $\pi$ between $G_1$ and $G_2$ is never revealed.

%==============================================================================
\section{Bonus Q3: Exploiting the Biased Verifier}
%==============================================================================

\subsection{Problem Overview}

In this section, we demonstrate that the soundness of the 3-coloring interactive proof relies critically on the \textbf{uniform randomness} of the Verifier's edge selection. The provided verifier has a bias: it preferentially challenges edges among vertices with smaller indices.

Our goal is to construct colorings that are \textit{invalid} (have conflicts) but concentrate those conflicts in rarely-challenged regions of the graph.

\subsection{Methodology}

\subsubsection{Graph Generation}
We generated 20 non-3-colorable graphs of varying sizes using the provided generator (treated as a black box, as required by the assignment). Each graph file stores $n$, $m$, and then $m$ edges.

We used sizes up to 1000 vertices, as suggested:
10, 15, 20, 30, 40, 50, 60, 70, 80, 100, 120, 150, 200, 250, 300, 400, 500, 600, 800, 1000.

\subsubsection{Adversarial Coloring Strategy}
The verifier is \emph{not} uniform: in our provided implementation (\texttt{mock\_zk\_verifier.py}), it samples from a \emph{biased subset} of edges with high probability (80\%), where both endpoints lie in the first $\lfloor n/3 \rfloor$ vertices. Therefore, a natural adversarial approach is to try to produce colorings that:
\begin{itemize}
    \item have few or no conflicts among the vertices that are likely to be checked, and
    \item ``push'' remaining conflicts to edges that are less likely to be sampled.
\end{itemize}

In our submission we do \textbf{not} implement a cheating prover; instead, we generate a fixed coloring file and use the provided \texttt{honest\_prover.py}, exactly as required.

\subsubsection{Coloring Algorithm (Implemented Heuristic)}
Our \texttt{solve\_q3.py} implements a simple local-search heuristic to \emph{minimize the total number of conflicts} (not necessarily to eliminate them, since the graphs are generated to be non-3-colorable). At a high level:
\begin{itemize}
    \item Start from a random 3-color assignment.
    \item Repeatedly pick a random vertex and attempt a recoloring step.
    \item Accept a recoloring greedily if it does not increase the number of conflicts incident to that vertex.
    \item Track the best coloring seen and output it.
\end{itemize}

The key point is that even a modest reduction in conflicts can yield a noticeably higher acceptance rate when the verifier's sampling distribution avoids many of the remaining conflicting edges.

\subsection{Experimental Results}

We ran each graph-coloring pair through 20 trials with 10 rounds per trial.

\begin{table}[h]
\centering
\begin{tabular}{@{}cccc@{}}
\toprule
\textbf{Graph ID} & \textbf{Size ($n$)} & \textbf{Conflicts} & \textbf{Acceptance Rate (\%)} \\ \midrule
1  & 10   & 6   & 45 \\
2  & 15   & 9   & 45 \\
3  & 20   & 7   & 70 \\
4  & 30   & 14  & 55 \\
5  & 40   & 15  & 35 \\
6  & 50   & 21  & 35 \\
7  & 60   & 17  & 70 \\
8  & 70   & 21  & 60 \\
9  & 80   & 29  & 60 \\
10 & 100  & 40  & 70 \\
11 & 120  & 24  & 75 \\
12 & 150  & 53  & 45 \\
13 & 200  & 64  & 65 \\
14 & 250  & 80  & 60 \\
15 & 300  & 94  & 75 \\
16 & 400  & 123 & 70 \\
17 & 500  & 154 & 55 \\
18 & 600  & 199 & 55 \\
19 & 800  & 258 & 45 \\
20 & 1000 & 302 & 70 \\ \bottomrule
\end{tabular}
\caption{Measured conflicts and acceptance rates (20 trials per graph, 10 rounds per trial). Despite all colorings being invalid (having conflicts), the biased verifier accepts them 35--75\% of the time.}
\label{tab:results}
\end{table}

\subsection{Analysis}

\subsubsection{Key Observations}
\begin{enumerate}
    \item \textbf{Size effects:} In practice, acceptance rates can change significantly with graph size, because both the number of edges and the verifier's biased sampling region scale with $n$.
    \begin{itemize}
        \item If many conflicts occur outside the verifier's frequently-sampled edge subset, they may be detected rarely.
        \item The same number of conflicts can be ``harder to find'' when sampling is biased away from where conflicts are.
    \end{itemize}
    
    \item \textbf{Bias exploitation:} With 80\% of challenges drawn from a biased subset of edges:
    \begin{itemize}
        \item Conflicts concentrated outside this subset are less likely to be detected.
        \item Even if the total number of conflicts is nonzero, the verifier may miss them for many rounds.
    \end{itemize}
    
    \item \textbf{Probability analysis:} Let $q_b$ be the probability that a \emph{biased} challenge edge is conflicting under our produced coloring, and let $q_u$ be the probability that an \emph{unbiased} (fully random) challenge edge is conflicting. Then the per-round rejection probability is:
    \[
    q = 0.8\,q_b + 0.2\,q_u,
    \]
    so the acceptance probability over $k$ rounds is:
    \[
    \Pr[\text{Accept}] = (1 - q)^k.
    \]
    This highlights why lowering $q_b$ (conflicts inside the biased subset) can disproportionately increase the overall acceptance probability.
\end{enumerate}

\subsubsection{Why This Doesn't Contradict Theory}

The theoretical soundness guarantee assumes:
\begin{enumerate}
    \item \textbf{Uniform Edge Selection:} The Verifier chooses edges uniformly at random.
    \item \textbf{Cryptographic Commitments:} The Prover commits to colors before seeing the challenge.
\end{enumerate}

Our experiment violates the first assumption. With uniform selection, each conflict edge has probability $1/m$ of being caught per round. With biased selection:
\[
\Pr[\text{catch conflict}] = 0.8 \cdot \Pr[\text{conflict in biased region}] + 0.2 \cdot \Pr[\text{conflict overall}]
\]

If we ensure few (or zero) conflicts in the biased region, the dominant term $0.8\,q_b$ becomes small, and the verifier's ability to catch conflicts relies mostly on the remaining 20\% of challenges.

\textbf{Lesson:} The randomness quality in the Verifier's challenges is essential for soundness. A predictable or biased Verifier can be exploited by an adversarial input distribution.

%==============================================================================
\section{Bonus Q4: Constructing Zero-Knowledge Proofs for NP}
%==============================================================================

\subsection{General Methodology}

Any problem in NP can be proven in zero-knowledge using the following construction:

\subsubsection{Step 1: Reduction to NP-Complete Problem}
Since 3-Coloring (or SAT, Hamiltonian Path, etc.) is NP-complete:
\begin{itemize}
    \item For any language $L \in \text{NP}$, there exists a polynomial-time reduction $f$ such that $x \in L \iff f(x) \in \textsc{3Color}$.
    \item The reduction preserves witness structure: a certificate for $x$ can be transformed into a valid 3-coloring.
\end{itemize}

\subsubsection{Step 2: Witness Translation}
\begin{itemize}
    \item The Prover holds a witness $w$ for $x \in L$.
    \item Using the reduction's logic, the Prover constructs the corresponding witness for the NP-complete problem (e.g., a valid 3-coloring of $f(x)$).
\end{itemize}

\subsubsection{Step 3: Apply Standard ZK Protocol}
\begin{itemize}
    \item Run the zero-knowledge protocol for the NP-complete problem.
    \item The Verifier is convinced that $f(x) \in \textsc{3Color}$, hence $x \in L$.
\end{itemize}

\subsubsection{Step 4: Cryptographic Commitments}
Commitments ensure:
\begin{itemize}
    \item \textbf{Hiding:} The Verifier cannot determine the committed value before opening.
    \item \textbf{Binding:} The Prover cannot change the committed value after committing.
\end{itemize}

Common schemes include Pedersen commitments (computationally binding, perfectly hiding) or hash-based commitments.

\subsubsection{Step 5: Probabilistic Soundness}
\begin{itemize}
    \item Each round catches a cheating Prover with probability at least $1/\text{poly}(n)$.
    \item After $k$ rounds, cheating probability drops to $(1 - 1/\text{poly}(n))^k$.
    \item For $k = \Omega(n \cdot \text{poly}(n))$, this becomes negligible.
\end{itemize}

%==============================================================================
\section{Bonus Q5: Interactive Proof for Hamiltonian Path}
%==============================================================================

\subsection{Problem Definition}
\[
\textsc{HP} = \{ G \mid G \text{ contains a Hamiltonian path} \}
\]

A Hamiltonian path visits every vertex exactly once.

\subsection{Zero-Knowledge Protocol}

\textbf{Setup:}
\begin{itemize}
    \item \textbf{Common Input:} Graph $G = (V, E)$ with $n$ vertices.
    \item \textbf{Prover's Secret:} A Hamiltonian path $P = (v_1, v_2, \ldots, v_n)$.
\end{itemize}

\textbf{Protocol (Single Round):}
\begin{enumerate}
    \item \textbf{Commitment:}
    \begin{itemize}
        \item Prover selects a random permutation $\pi$ of vertices.
        \item Prover computes $H = \pi(G)$ and commits to its adjacency matrix entry-by-entry.
    \end{itemize}
    
    \item \textbf{Challenge:}
    \begin{itemize}
        \item Verifier sends random bit $b \in \{1, 2\}$.
    \end{itemize}
    
    \item \textbf{Response:}
    \begin{itemize}
        \item If $b = 1$: Prover reveals $\pi$ and opens all adjacency matrix entries.
        \item If $b = 2$: Prover reveals the path $\pi(P)$ in $H$ by opening only the $n-1$ path edges.
    \end{itemize}
    
    \item \textbf{Verification:}
    \begin{itemize}
        \item If $b = 1$: Verifier checks $H = \pi(G)$.
        \item If $b = 2$: Verifier checks that the revealed edges form a valid Hamiltonian path in $H$.
    \end{itemize}
\end{enumerate}

\subsection{Analysis}

\subsubsection{Completeness}
If $G$ has a Hamiltonian path $P$:
\begin{itemize}
    \item For $b = 1$: Prover reveals the isomorphism. Check passes trivially.
    \item For $b = 2$: Prover reveals $\pi(P)$, which is a valid Hamiltonian path in $H = \pi(G)$.
\end{itemize}

\subsubsection{Soundness}
If $G$ has no Hamiltonian path:
\begin{itemize}
    \item To pass $b = 1$: Need $H \cong G$, so $H$ also has no Hamiltonian path.
    \item To pass $b = 2$: Need $H$ to have a Hamiltonian path.
\end{itemize}

These requirements are contradictory. The Prover can satisfy at most one, so:
\[
\Pr[\text{success per round}] \leq \frac{1}{2}
\]

After $k$ rounds: $\Pr[\text{cheat}] \leq (1/2)^k$.

\subsubsection{Zero-Knowledge}
\begin{itemize}
    \item For $b = 1$: Verifier sees a random isomorphic copy of $G$ (no information about the path).
    \item For $b = 2$: Verifier sees $n-1$ edges forming a path on random vertex labels. Since $\pi$ is random, the path vertices are a random permutation of $\{1, \ldots, n\}$---this is just a random sequence with no correlation to the original path $P$.
\end{itemize}

\textbf{Simulator:} 
\begin{enumerate}
    \item Pick random $b'$.
    \item If $b' = 1$: Generate random $H = \sigma(G)$, output $\sigma$.
    \item If $b' = 2$: Create $H$ with a random Hamiltonian path embedded.
    \item If $b \neq b'$: Rewind and retry.
\end{enumerate}

%==============================================================================
\section{Conclusion}
%==============================================================================

This project demonstrated the elegant interplay between randomness, computational complexity, and cryptography in interactive proof systems. Key takeaways include:

\begin{enumerate}
    \item \textbf{Zero-Knowledge is Powerful:} NP-complete problems can be proven without revealing witnesses, enabling applications in authentication and secure computation.
    
    \item \textbf{Randomness is Essential:} The Verifier's random challenges prevent the Prover from ``gaming'' the system. Biased randomness (as in Q3) breaks soundness.
    
    \item \textbf{Repetition Amplifies Security:} While single-round soundness may be weak ($1 - 1/m$ or $1/2$), polynomial repetition achieves exponentially small error.
    
    \item \textbf{Assumptions Matter:} The simplified protocol in Q3 omits cryptographic commitments and also uses a biased sampling strategy. In the full theoretical protocol, commitments (binding/hiding) and uniform randomness are key ingredients in the soundness and zero-knowledge arguments.
\end{enumerate}

%==============================================================================
\section*{References}
%==============================================================================
\begin{enumerate}
    \item Goldreich, O., Micali, S., \& Wigderson, A. (1991). Proofs that yield nothing but their validity. \textit{Journal of the ACM}.
    \item Goldwasser, S., Micali, S., \& Rackoff, C. (1989). The knowledge complexity of interactive proof systems. \textit{SIAM Journal on Computing}.
    \item Arora, S., \& Barak, B. (2009). \textit{Computational Complexity: A Modern Approach}. Cambridge University Press.
\end{enumerate}

\end{document}
